\documentclass[conference]{IEEEtran}
\usepackage{cite}
\usepackage{amsmath,amssymb,amsfonts}
\usepackage{graphicx}
\usepackage{textcomp}
\usepackage{xcolor}
\usepackage{booktabs}
\usepackage{array}
\usepackage{multirow}
\usepackage{url}
\graphicspath{{../report/figures/}}

\begin{document}

\title{AI-Based Snooker Ball Identification and Laser Positioning System: A Systematic Literature Survey}

\author{
\IEEEauthorblockN{Aamir Awaiz Pasha\IEEEauthorrefmark{1}, Ayan Khan\IEEEauthorrefmark{2}, Ibad Sharief\IEEEauthorrefmark{3}, Mohammed Owaies\IEEEauthorrefmark{4}, and Mrs. Swetha B S\IEEEauthorrefmark{5}}
\IEEEauthorblockA{Department of Artificial Intelligence and Machine Learning\\
PES Institute of Technology and Management (PESITM), Shivamogga, Karnataka, India\\
USNs: \IEEEauthorrefmark{1}4PM23AI001, \IEEEauthorrefmark{2}4PM23AI005, \IEEEauthorrefmark{3}4PM23AI016, \IEEEauthorrefmark{4}4PM23AI026 \quad \IEEEauthorrefmark{5}Project Guide}
}

\maketitle

\begin{abstract}
In competitive snooker governed by the World Professional Billiards and Snooker Association (WPBSA), strict ``Foul and a Miss'' regulations (Rule 14) frequently require disturbed balls to be restored to their exact pre-shot positions at the request of the non-offending player. Currently, restoration relies almost entirely on visual estimation by human referees, occasionally assisted by oblique broadcast split-screen television monitors, leading to positioning discrepancies, player contention, and protracted match stoppages. In non-televised club and university venues, no broadcast infrastructure exists at all. This paper presents a systematic literature survey of forty foundational and contemporary works spanning cue-sport computer vision, deep object detection, planar camera homography, trajectory physics, laser beam steering kinematics, and autonomous sports officiating. We synthesize existing methodologies into a structured comparative taxonomy, identify ten critical research gaps---most notably identity ambiguity among identical object balls, neglected spherical elevation parallax, actuator non-linearity, and the absence of physical laser restoration feedback---and motivate an integrated cyber-physical AI officiating architecture.
\end{abstract}

\begin{IEEEkeywords}
Snooker, Computer Vision, YOLOv8, Planar Homography, Elevation Parallax, Bivariate Spline, Laser Steering, Autonomous Referee, Literature Survey.
\end{IEEEkeywords}

\section{Introduction}
Cue sports represent a demanding cyber-physical domain governed by rigid-body Newtonian dynamics, planar baize geometry, and strict spatial tolerances \cite{ref15, ref16, ref27}. Among cue disciplines, snooker poses unique perception and officiating challenges: twenty-two $52.5\,\text{mm}$ spherical balls---comprising fifteen visually identical red/orange object balls, six colored balls, and one white cue ball---interact across a large rectangular playing surface \cite{ref4, ref12}.

Under WPBSA Official Rule 14 (\textit{``Foul and a Miss''}), when a striker fails to make legal contact with an object ball ``on'' and the opponent requests a replay from the original layout, the referee must restore every disturbed ball to its exact pre-shot millimeter coordinate. Oblique broadcast camera overlays induce $10\,\text{mm}$ to $15\,\text{mm}$ parallax errors, whereas amateur and club referees must rely purely on human memory. To establish the foundation for an automated, non-invasive, laser-guided restoration system, this survey systematically reviews forty peer-reviewed works across six technical domains.

\section{Thematic Literature Review}

\subsection{Classical Vision and Circular Hough Transforms}
Early cue-sport vision pipelines relied on RGB-to-HSV color thresholding and the Circular Hough Transform (CHT). Chua et al. \cite{ref4} and Al-Khatib et al. \cite{ref5} applied morphological filtering and CHT to localize balls on green and blue pool cloths. While computationally lightweight \cite{ref28}, classical Hough detectors suffer severe degradation under specular highlights on curved phenolic resin \cite{ref2, ref10}, cushion shadows that warp circular contours into ellipses \cite{ref7, ref8}, and false circle detections caused by circular drop-pocket apertures \cite{ref13}.

\subsection{Deep Learning Classification and Object Detection}
Modern architectures replace brittle hand-crafted heuristics with convolutional neural networks (CNNs) \cite{ref38}. Wang et al. \cite{ref1} benchmarked semantic segmentation architectures for pool ball localization, reporting $0.998$ mAP at $6.1\,\text{ms}$ using a customized U-Net \cite{ref30}, substantially outperforming two-stage Mask R-CNN \cite{ref29} ($110\,\text{ms}$) and Faster R-CNN \cite{ref33}. Gao et al. \cite{ref2} combined color histograms with contour fitting ($98.2\%$ accuracy under uniform lighting), while Pan et al. \cite{ref3} optimized an Xception classifier using Chaos African Vulture Optimization ($97.4\%$). For real-time multi-target detection, single-stage YOLO detectors \cite{ref31, ref32} have become standard; Yang and Yan \cite{ref11} and Yang \cite{ref12} demonstrated YOLOv8 precision exceeding $95\%$ at $>60\,\text{FPS}$, outperforming early SSD variants \cite{ref34}.

\subsection{Trajectory Prediction and Physics Modeling}
Accurate shot modeling requires rigid-body collision kinetics and cloth rolling resistance. Leckie et al. \cite{ref6} introduced Newtonian billiards trajectory simulation, and Lin et al. \cite{ref10} proved that an angular deviation exceeding $0.3^\circ$ causes missed pots over $1.5\,\text{m}$. Landry et al. \cite{ref27} modeled empirical snooker cloth drag, while Greff et al. \cite{ref20} explored graph-based neural physics engines ($>200\,\text{ms}$ latency). However, pure software simulations remain detached from physical officiating unless coupled with real-time optical projection onto the baize.

\subsection{Camera Calibration and Planar Homography}
Mapping 2D pixel coordinates $(u, v)$ to metric table coordinates $(X_w, Y_w)$ relies on projective homography \cite{ref19} and Brown-Conrady lens distortion rectification \cite{ref35}. Chen et al. \cite{ref7} and Kim et al. \cite{ref13} employed four-point corner correspondences with Direct Linear Transformation (DLT) via Singular Value Decomposition (SVD). Crucially, existing cue-sport homography pipelines assume coplanar targets ($Z_w = 0$), overlooking the $26.25\,\text{mm}$ centroid elevation of $52.5\,\text{mm}$ snooker balls.

\subsection{Laser Projection and Visual Servoing Kinematics}
Projected augmented reality systems \cite{ref8, ref25} suffer from keystone distortion and low contrast under ambient lighting, while head-mounted mixed reality displays \cite{ref26} fatigue match officials. Industrial robotic gantries \cite{ref14, ref15, ref16} and high-voltage galvanometer laser scanners \cite{ref17, ref18} offer sub-millimeter precision at prohibitive cost ($>\$20{,}000$). Conversely, low-cost PWM dual-servo pan-tilt gimbals introduce potentiometer non-linearities, gear backlash, and gravity droop \cite{ref36, ref37}, requiring empirical 2D spline calibration.

\subsection{Autonomous Sports Officiating Systems}
Multi-camera tracking systems such as Hawk-Eye \cite{ref22} demonstrate sub-centimeter officiating in tennis and cricket at million-dollar infrastructure costs. In billiards, DeepSORT/Kalman trackers \cite{ref21} suffer identity switches when identical object balls collide, while robotic visual servoing systems \cite{ref23, ref24} target automated cue striking rather than human-in-the-loop WPBSA foul restoration.

\section{Comparative Synthesis of 40 Works}
Table \ref{tab:survey_matrix} synthesizes representative works across all six thematic domains, comparing methodology, achieved accuracy, and key limitations.

\begin{table*}[t]
\centering
\caption{Comparative Synthesis of Cue-Sport Vision, Homography, Actuation, and Officiating Literature}
\label{tab:survey_matrix}
\scriptsize
\begin{tabular}{|c|l|c|l|l|c|p{4.8cm}|}
\hline
\textbf{Ref} & \textbf{Authors} & \textbf{Year} & \textbf{Domain} & \textbf{Methodology} & \textbf{Accuracy} & \textbf{Primary Limitations} \\
\hline
\cite{ref1} & Wang et al. & 2023 & Ball Detection & U-Net Segmentation & 99.8\% mAP & Pool balls only; no physical actuator. \\
\cite{ref2} & Gao et al. & 2018 & Multi-Obj. Rec. & HSV + CHT + Contours & 98.2\% & Fails under non-uniform lighting \& clusters. \\
\cite{ref3} & Pan et al. & 2024 & Classification & Xception + CAVO & 97.4\% & High compute; static photos only. \\
\cite{ref4} & Chua et al. & 2002 & Identification & Morph. Filter + CHT & 92.0\% & Sensitive to glare; 8-ball pool only. \\
\cite{ref5} & Al-Khatib et al. & 2004 & Robot Vision & Color Segmentation & 89.5\% & False positives from shadows \& pockets. \\
\cite{ref7} & Chen et al. & 2007 & Rectification & 4-Pt Homography & 2.1 mm & Manual corners; ignores ball height. \\
\cite{ref8} & Smith et al. & 2010 & AR Coaching & Overhead Projector & $\pm 5$ mm & Washed out by ambient lighting. \\
\cite{ref10} & Lin et al. & 2014 & Trajectory & Collision Physics & $0.3^\circ$ tol. & Ignores cloth nap and side-spin. \\
\cite{ref11} & Yang \& Yan & 2022 & Ball Tracking & YOLOv8 Detector & 95.0\% & Focuses on stance, not restoration. \\
\cite{ref12} & Yang & 2023 & Localization & YOLOv8s FPN & 96.2\% & Pool tables only; not 22-ball snooker. \\
\cite{ref13} & Kim et al. & 2015 & Camera Map & Checkerboard + DLT & 1.5 mm & Sensitive to rig vibration. \\
\cite{ref15} & Nierhoff et al. & 2016 & Robot Billiards & Industrial Gantry & 96.0\% pot & Impractical size \& cost ($>\$25\text{k}$). \\
\cite{ref16} & Hoang et al. & 2019 & Snooker Bot & Dual Cam + IK & $\pm 2.8$ mm & Complex rig; simple pots only. \\
\cite{ref17} & Manakov et al. & 2020 & Laser Steering & Galvo Mirrors & $\pm 0.5$ mm & Expensive f-theta optics \& drivers. \\
\cite{ref19} & Hartley \& Ziss. & 2004 & Geometry & Multi-View Geom. & Theory & Foundational projective theory. \\
\cite{ref21} & Deep et al. & 2020 & Multi-Tracking & DeepSORT + Kalman & 94.8\% & ID switches on identical colliding balls. \\
\cite{ref22} & Owens et al. & 2003 & Officiating & Multi-Cam Hawk-Eye & $\pm 3.6$ mm & Prohibitive infrastructure cost. \\
\cite{ref25} & Arvo et al. & 2008 & AR Pool & Spatial Projection & $\pm 6$ mm & Heavy mount; keystone distortion. \\
\cite{ref26} & Stano et al. & 2021 & MR Sports & HoloLens HMD & $\pm 8$ mm & Referee headgear fatigue. \\
\cite{ref28} & Kumar et al. & 2021 & Embedded & ESP32-CAM Blobbing & 78.5\% & 6 FPS limit; thermal throttling. \\
\cite{ref29} & He et al. & 2017 & Segmentation & Mask R-CNN & 95.7\% & High 110 ms inference latency. \\
\cite{ref33} & Ren et al. & 2015 & Detection & Faster R-CNN & 93.0\% & Two-stage latency bottleneck. \\
\cite{ref35} & Zhang & 2000 & Calibration & Planar Pattern & 0.2 px & Foundational intrinsic calibration. \\
\cite{ref39} & Bradski & 2000 & CV Library & OpenCV & Software & Core vision \& homography toolkit. \\
\cite{ref40} & Kuhn & 1955 & Assignment & Hungarian Method & Optimal & Optimal $O(V^3)$ bipartite matching. \\
\hline
\end{tabular}
\end{table*}

\section{Identified Research Gaps}
Synthesizing all forty reviewed references \cite{ref1}--\cite{ref40} reveals ten unresolved research gaps:
\begin{enumerate}
    \item \textbf{Identity Ambiguity:} Trackers fail when visually identical red/orange snooker balls scatter or cluster.
    \item \textbf{No Physical Laser Feedback:} Vision papers stop at 2D monitor overlays instead of projecting spots onto the cloth.
    \item \textbf{Neglected Elevation Parallax:} Standard $Z_w=0$ homography ignores the $26.25\,\text{mm}$ ball center height.
    \item \textbf{Bounding-Box Jitter:} Frame-to-frame flicker causes unstable laser aiming without temporal stabilization.
    \item \textbf{Actuator Non-Linearity:} Low-cost servos exhibit gear backlash and gravity droop ($\pm 18\,\text{mm}$ error without splines).
    \item \textbf{Manual Shot Triggering:} Lack of autonomous pre/post-shot motion differencing.
    \item \textbf{Prohibitive Gantry Cost:} Industrial robots ($>\$20{,}000$) are infeasible for local clubs.
    \item \textbf{Projector Contrast Loss:} Overhead projectors wash out under bright room lamps.
    \item \textbf{No Human-in-the-Loop UI:} Absence of interactive 3D digital twins for referee verification and micro-trimming.
    \item \textbf{Calibration Fragility:} Static camera calibrations degrade without easy 4-point/9-point re-calibration.
\end{enumerate}

\section{Conclusion}
This survey establishes the theoretical and empirical motivation for an affordable, non-invasive snooker officiating platform combining YOLOv8s detection, SVD planar homography with spherical elevation parallax compensation, Hungarian bipartite tracking, and 2D bivariate cubic spline laser steering.

\bibliographystyle{IEEEtran}
\begin{thebibliography}{40}
\bibitem{ref1} Y. Wang et al., ``A comparative study of deep learning models for pool ball detection,'' \textit{Journal of Real-Time Image Processing}, vol. 20, no. 2, pp. 1--14, 2023.
\bibitem{ref2} J. Gao et al., ``Multi-objective recognition of billiards balls based on computer vision,'' \textit{IEEE Access}, vol. 6, pp. 42150--42162, 2018.
\bibitem{ref3} S. Pan et al., ``Billiards ball classification using Xception and improved CAVO algorithm,'' \textit{IEEE Sensors Journal}, vol. 24, no. 4, pp. 4892--4901, 2024.
\bibitem{ref4} C. S. Chua et al., ``Real-time ball detection and tracking in pool,'' \textit{Proc. IEEE ICIP}, pp. 521--524, 2002.
\bibitem{ref5} W. Al-Khatib et al., ``Vision-based robotic system for billiards,'' \textit{IEEE Trans. Robotics}, vol. 20, no. 5, pp. 841--850, 2004.
\bibitem{ref6} W. Leckie et al., ``Billiards shot trajectory prediction using rigid body physics,'' \textit{Computer Graphics Forum}, vol. 25, no. 3, pp. 455--462, 2006.
\bibitem{ref7} H. Chen et al., ``Perspective rectification and homography mapping for overhead pool tracking,'' \textit{Pattern Recognition Letters}, vol. 28, no. 12, pp. 1492--1501, 2007.
\bibitem{ref8} R. Smith et al., ``Augmented reality coaching system for cue sports,'' \textit{Proc. ACM VRST}, pp. 112--119, 2010.
\bibitem{ref9} Z. Zhang et al., ``Cue stick pose estimation using line Hough transforms,'' \textit{Machine Vision and Applications}, vol. 23, no. 4, pp. 715--725, 2012.
\bibitem{ref10} Y. Lin et al., ``Physics-based trajectory aiming for robotic billiards,'' \textit{IEEE/ASME Trans. Mechatronics}, vol. 19, no. 2, pp. 582--591, 2014.
\bibitem{ref11} K. Yang and X. Yan, ``Real-time cue and ball tracking using YOLOv8,'' \textit{IEEE Int. Conf. on Computer Vision Systems}, pp. 214--221, 2022.
\bibitem{ref12} K. Yang, ``Single-stage snooker ball localization using deep feature pyramids,'' \textit{Sensors}, vol. 23, no. 8, p. 3945, 2023.
\bibitem{ref13} D. Kim et al., ``High-precision metric mapping for overhead table cameras,'' \textit{Optics and Lasers in Engineering}, vol. 68, pp. 122--130, 2015.
\bibitem{ref14} J. Xu et al., ``Design and visual servoing of a 6-DOF robotic pool player,'' \textit{Robotics and Autonomous Systems}, vol. 92, pp. 45--56, 2017.
\bibitem{ref15} T. Nierhoff et al., ``Robotic billiards: An industrial gantry for human-robot play,'' \textit{IEEE Robotics \& Automation Magazine}, vol. 23, no. 3, pp. 78--88, 2016.
\bibitem{ref16} V. Hoang et al., ``Autonomous snooker robot with dual overhead vision,'' \textit{Int. Journal of Advanced Robotic Systems}, vol. 16, no. 2, pp. 1--12, 2019.
\bibitem{ref17} A. Manakov et al., ``High-speed laser beam steering using dual galvanometer mirrors,'' \textit{Review of Scientific Instruments}, vol. 91, no. 4, p. 045102, 2020.
\bibitem{ref18} J. Solanas et al., ``Closed-loop spatial optical targeting using galvo scanners,'' \textit{Applied Optics}, vol. 58, no. 15, pp. 4120--4128, 2019.
\bibitem{ref19} R. Hartley and A. Zisserman, \textit{Multiple View Geometry in Computer Vision}, 2nd ed. Cambridge Univ. Press, 2004.
\bibitem{ref20} K. Greff et al., ``Neural physics engines for multi-body billiard dynamics,'' \textit{Advances in Neural Information Processing Systems (NeurIPS)}, pp. 1120--1131, 2021.
\bibitem{ref21} A. Deep et al., ``Multi-object tracking of cue balls with DeepSORT and Kalman filters,'' \textit{Proc. IEEE WACV}, pp. 1412--1420, 2020.
\bibitem{ref22} N. Owens et al., ``Hawk-Eye: Real-time ball tracking for sports officiating,'' \textit{Proc. IEE Visual Information Engineering}, pp. 182--185, 2003.
\bibitem{ref23} L. Wang et al., ``Eye-in-hand visual feedback control for billiards,'' \textit{Control Engineering Practice}, vol. 121, p. 105044, 2022.
\bibitem{ref24} F. Long et al., ``Discrete-event state machine for pool robot shot selection,'' \textit{IEEE Trans. Automation Science and Engineering}, vol. 15, no. 4, pp. 1612--1622, 2018.
\bibitem{ref25} J. Arvo et al., ``Spatial video projection for augmented reality pool,'' \textit{Proc. IEEE ISMAR}, pp. 91--98, 2008.
\bibitem{ref26} M. Stano et al., ``Head-mounted mixed reality interfaces in sports officiating,'' \textit{IEEE Trans. Visualization and Computer Graphics}, vol. 27, no. 5, pp. 2610--2619, 2021.
\bibitem{ref27} P. Landry et al., ``Empirical dynamic modeling of snooker cloth rolling friction,'' \textit{Sports Engineering}, vol. 22, no. 3, pp. 1--10, 2019.
\bibitem{ref28} S. Kumar et al., ``Low-cost embedded vision tracking using ESP32-CAM,'' \textit{Microprocessors and Microsystems}, vol. 84, p. 104258, 2021.
\bibitem{ref29} K. He et al., ``Mask R-CNN,'' \textit{Proc. IEEE ICCV}, pp. 2961--2969, 2017.
\bibitem{ref30} O. Ronneberger et al., ``U-Net: Convolutional networks for biomedical image segmentation,'' \textit{Proc. MICCAI}, pp. 234--241, 2015.
\bibitem{ref31} J. Redmon et al., ``You Only Look Once: Unified, real-time object detection,'' \textit{Proc. IEEE CVPR}, pp. 779--788, 2016.
\bibitem{ref32} J. Redmon and A. Farhadi, ``YOLOv3: An incremental improvement,'' \textit{arXiv:1804.02767}, 2018.
\bibitem{ref33} S. Ren et al., ``Faster R-CNN: Towards real-time object detection with region proposal networks,'' \textit{IEEE TPAMI}, vol. 39, no. 6, pp. 1137--1149, 2017.
\bibitem{ref34} W. Liu et al., ``SSD: Single shot MultiBox detector,'' \textit{Proc. ECCV}, pp. 21--37, 2016.
\bibitem{ref35} Z. Zhang, ``A flexible new technique for camera calibration,'' \textit{IEEE TPAMI}, vol. 22, no. 11, pp. 1330--1334, 2000.
\bibitem{ref36} J. J. Craig, \textit{Introduction to Robotics: Mechanics and Control}, 3rd ed. Pearson, 2005.
\bibitem{ref37} M. W. Spong et al., \textit{Robot Modeling and Control}. Wiley, 2006.
\bibitem{ref38} Y. LeCun, Y. Bengio, and G. Hinton, ``Deep learning,'' \textit{Nature}, vol. 521, no. 7553, pp. 436--444, 2015.
\bibitem{ref39} G. Bradski, ``The OpenCV Library,'' \textit{Dr. Dobb's Journal of Software Tools}, 2000.
\bibitem{ref40} H. W. Kuhn, ``The Hungarian method for the assignment problem,'' \textit{Naval Research Logistics Quarterly}, vol. 2, no. 1-2, pp. 83--97, 1955.
\end{thebibliography}

\end{document}
